\section{Domain Analysis}
\label{sec:domainanalysis}

The domain is the deployment configuration of a small production estate: one
Kubernetes cluster of seven nodes, ten namespaces, and roughly thirty
long-running processes belonging to about twenty applications. The estate is
real and its current configuration is hand-maintained, which is what makes it
usable as evidence: every concept below was extracted from configuration that
exists, and several modelling decisions are justified by a defect the existing
configuration already has.

\subsection{Where the Repetition Is}

Four observations from the existing configuration drive the whole language.

\textbf{One fact is written down many times.} The hostname of one application
appears in seven authoritative places across three repositories, and two
conformance tests exist for no purpose other than detecting when the seven
disagree. A port appears in the container specification, in the service, in the
probe and in the monitoring rule. The guard was cheaper to write than the fix,
which is the situation a generated deployment set removes.

\textbf{One derivation is performed many times by hand.} Four deployments carry
the same rollout block, arrived at independently, with the reasoning recorded in
comments that no tool can read. Comments record that a recreate strategy opened
a zero-pod window, and that a JVM cold start needs a startup budget an order of
magnitude larger than a static bundle's. That is a rule about a class of
process, expressed four times as a constant.

\textbf{Values are stated in the wrong vocabulary.} The existing configuration
declares mechanisms, not requirements: a rollout strategy rather than whether
the application may be taken down during a cutover, a node label rather than
what the process needs from a machine. A process whose author forgets to declare
a recreate strategy over a volume that cannot be attached twice appears to work
on a one-node cluster and fails the first time a second node exists.

\textbf{Some values cannot be written by one application at all.} A hostname
must be unique across the estate; a graphics card is held by one process at a
time; the edge is a shared, finite resource. These values are contended, and no
application can decide them from inside its own repository.

\subsection{The Concepts the Domain Requires}

The analysis produces two groups of concepts, and the split between them is the
main modelling decision of the project (Section~\ref{sec:decisions}).

The \emph{structural} concepts are the ones a reader would expect: a project
groups applications, an application groups the processes released together,
a process offers named ports and depends on the ports of others, and an
application is reachable at a hostname whose paths route to its processes.

The \emph{operational} concepts are the ones that carry most of the generated
output: what a process needs from a machine, what its data is worth if it is
lost, which secrets it may read and how it receives them, what answers a health
check, how long a cold start may legitimately take, and whether a new version
may run beside the old one during a cutover. These have no counterpart in a
structural architecture notation, and they are where the interesting derivations
live.

\subsection{Notation: Why the Metamodels Are Drawn as Class Diagrams}
\label{sec:notation}

The Task 0 feedback asked why UML class diagrams are used rather than a
structural notation such as a deployment or component diagram. The two answer
different questions, and this project needs both answers.

A metamodel defines the \emph{abstract syntax} of a language: which concepts
exist, which properties they carry, and how they may be connected. Ecore, the
metamodelling language used here, is an implementation of EMOF, and EMOF is
itself defined in terms of classes, attributes, references and multiplicities.
A class diagram is therefore not one of several ways to picture an Ecore
package; it is the notation whose constructs correspond one-to-one to the
constructs the metamodel is built from. A deployment diagram, by contrast, shows
\emph{one} arrangement of artifacts on nodes. It is an instance-level notation,
so it can illustrate a model of this language, but it cannot define the
language.

The distinction matters because this project has models at both levels. The
class diagrams in Sections~\ref{sec:sourcemetamodel}
and~\ref{sec:targetmetamodel} define the two languages. A deployment-style
picture of one estate would be a rendering of a \emph{model}, at the level the
example models in Section~\ref{sec:examplemodels} occupy, and that is where
ArchiMate and C4 become relevant.

\subsection{Relation to ArchiMate and the C4 Model}
\label{sec:archimate}

Both notations were examined, and the comparison is worth reporting because it
is a check on the vocabulary: if the language named its concepts differently
from the notations practitioners already use, that would be evidence the naming
was wrong.

\begin{table}[!htbp]
\centering
\small
\begin{tabular}{>{\raggedright\arraybackslash}p{0.22\linewidth}>{\raggedright\arraybackslash}p{0.34\linewidth}>{\raggedright\arraybackslash}p{0.34\linewidth}}
\hline
\textbf{This language} & \textbf{ArchiMate} & \textbf{C4 model} \\
\hline
Process & Application Component, deployed as an Artifact assigned to a Node &
  Container (a separately deployable, runnable unit) \\
Application & Application Collaboration, or a Grouping & no direct counterpart \\
Project & Grouping & Software System \\
Surface (a named port) & Application or Technology Interface & no first-class concept \\
Dependency edge & Serving relationship & a relationship between containers \\
Node contract & Node, with properties & Deployment Node \\
Exposure and Route & partially Path and Communication Network & shown as a relationship only \\
Placement dimension & no counterpart & no counterpart \\
Durability class & no counterpart & no counterpart \\
Secret grant & no counterpart & no counterpart \\
Probe, startup budget, cutover & no counterpart & no counterpart \\
\hline
\end{tabular}
\caption{The Source Language's Concepts Against ArchiMate and the C4 Model}
\label{tab:archimate}
\end{table}

Two conclusions follow from Table~\ref{tab:archimate}. The first is
reassuring: for every structural concept, both notations have an element with
essentially the same meaning, and the names agree. A process is a C4 container
and an ArchiMate application component deployed as an artifact; a dependency is
a serving relationship. Nothing in the structural half of the language is
invented.

The second conclusion is the reason neither notation is adopted as the source
language. The lower half of Table~\ref{tab:archimate} is empty on both sides.
ArchiMate and C4 describe how a system is arranged; they do not describe what an
application requires of the platform that runs it. Placement dimensions,
durability classes, secret grants, probes, startup budgets and cutover
requirements have no elements in either notation, and those concepts are
precisely what the generated files are derived from. Expressing them would mean
attaching free-form properties to ArchiMate elements, which gives up the one
thing a metamodel is for: a closed vocabulary that a constraint can be written
against. Section~\ref{sec:alternatives} records this as the rejected
alternative it is.

Where both notations remain useful is as a \emph{view} of a model rather than as
the model itself. An ArchiMate technology-layer view or a C4 deployment diagram
of an estate is derivable from a resolved model, because a resolved model
contains every fact such a view needs: which process runs where, which
application serves which hostname, and which process depends on which. That
derivation is a second model-to-text transformation over the same target
metamodel, and it is recorded as possible future work rather than promised
here, in keeping with the scope discipline of Section~\ref{sec:introduction}.
